%% ma: L12-B08-0136
%% lop: 12
%% chuong: 2
%% bai: 8
%% dang: 12.8.15
%% loai: TN4
%% muc: TH
%% dap_an: "A"
%% nguon: "CĐ2.3.pdf giáo viên gửi 27/10/2026 (M07, Lê Minh Khôi – Nguyễn Anh Khôi, Chương II Chủ đề 3 Biểu thức tọa độ của các phép toán vectơ), Câu 51"
%% kien_thuc: "Bài 8 (biểu thức tọa độ của phép toán vectơ), lớp 10 (hệ phương trình bậc nhất ba ẩn)"
%% trang_thai: chinh-thuc
\begin{ex}
Trong không gian $Oxyz$, cho các vectơ $\vec a=(1;2;1)$, $\vec b=(-2;3;4)$, $\vec c=(0;1;2)$ và $\vec d=(4;2;0)$. Biết rằng $\vec d=x\cdot\vec a+y\cdot\vec b+z\cdot\vec c$ với $x,y,z\in\mathbb R$. Tính giá trị của biểu thức $B=x+y+z$.
\choice{\True $B=2$}{$B=3$}{$B=4$}{$B=5$}
\loigiai{$\begin{cases}x-2y=4\\2x+3y+z=2\\x+4y+2z=0\end{cases}\Leftrightarrow x=2$, $y=-1$, $z=1$. Vậy $B=2$.}
\end{ex}
